\chapter{Discussion, Limitations, and Future Work}
\label{chap:discussion}

This chapter discusses the main findings from Experiment~1 and the
implications of the conflict-handling method introduced in Experiment~2. The results
show that explicit traceability rules improve master-prompt consolidation.
However, updating the specification after every prompt requires the agent to
make more decisions about how each new instruction should be interpreted,
which can introduce new errors. They also show that traceability alone is not sufficient when the
master prompt is handed to a later session, because new prompts may conflict
with already-settled requirements.

\section{Discussion}

\subsection{Traceability in Single-Pass Consolidation}

The main RQ1 finding is that a direct traceability rule is more effective than
a general request to preserve the meaning of the prompt history. The
\emph{strict} strategy performs better than \emph{basic} in all six
agent--project comparisons and better than \emph{loose} in five of them
(\cref{sec:results-18}).

However, the results do not show that adding more instructions always improves
the output. The \emph{loose} strategy occasionally performs worse than
\emph{basic}, while Codex performs slightly better with \emph{loose} than
with \emph{strict} on the simple project. Therefore, what matters is not the
length of the consolidation prompt, but whether it clearly limits the output
to requirements found in the original history.

Even the \emph{strict} strategy does not remove all non-traceable content.
This suggests that prompting can reduce unsupported additions, but cannot
guarantee a fully traceable master prompt.

\subsection{The Trade-Off in Spec-Guided Consolidation}
\label{sec:discussion-tradeoff}

The \emph{spec-guided} strategy improves on \emph{strict} in four of the six
agent--project comparisons. It performs well for all three agents on the
simple project, but only Codex improves on the hard project
(\cref{sec:rq2-results,tab:rq2-vs-strict}).

Maintaining the specification after every prompt reduces the need to
reconstruct the full project history at the end. However, it also requires the
agent to make a decision on every turn: whether the new prompt adds, changes,
removes, or conflicts with an existing requirement.

This becomes difficult when prompts contain conversational wording or several
changes to the same topic. Text that explains why a feature is requested may
be copied into the specification even though it is not itself a requirement.
Likewise, an instruction such as ``Change navigation buttons'' must be
rewritten as the resulting requirement rather than preserved as an editing
command.

The simple project contains short and direct prompts, so these interpretation
risks are limited. The hard project contains longer revision chains and more
conversational language, creating more opportunities for small errors to enter
the maintained specification.

The \emph{spec-guided} strategy should therefore be understood as an
alternative workflow, not as a guaranteed improvement over single-pass
consolidation. It reduces one large reconstruction problem, but replaces it
with many smaller interpretation decisions.

\subsection{Interpreting the Evaluation Scores}

The requirement-level evaluation helps identify which parts of a candidate
need inspection, but a low AlignScore does not always mean that the model
invented a requirement.

The manual review identified three main causes of low scores: the same
requirement expressed with different wording, conversational content that does
not belong in the final specification, and genuinely hallucinated details.
The untraceable rate combines these cases and should therefore be treated as a
screening measure rather than a direct hallucination rate.

The evaluation also measures only whether the requirements included in a
candidate are supported by the ground truth. It does not measure whether valid
requirements were omitted. A high score can therefore indicate strong
traceability without proving that the master prompt is complete.

\subsection{Semantic Conflicts After Project Handoff}

RQ3 considers how later prompts are handled after a settled master prompt is
handed to a new session. A requirement may conflict with an earlier decision
even when the two use different wording and do not create a direct text-level
clash.

The Experiment~2 governance instruction checks each later prompt against the
settled master prompt. Compatible additions are accepted, clear revisions
update the affected requirement, and unclear conflicts are listed under
\texttt{\#\# Unresolved conflicts} for the user to decide.

The method is demonstrated on both project histories using five constructed
conflicts and two compatible controls. Its purpose is to prevent an unclear
conflict from being silently treated as a project decision while allowing
compatible work to continue.

\section{Limitations and Threats to Validity}
\label{sec:discussion-limitations}

\begin{itemize}[leftmargin=*]

\item \textbf{Limited dataset.}
The evaluation uses two author-created web application histories, five
constructed semantic conflicts, and two compatible controls. The findings may
not generalize to larger projects, other domains, or different prompt-writing
styles.

\item \textbf{Limited repeated runs.}
Experiment~1 uses one scored candidate for each condition, while
Experiment~2 remains a small qualitative study. Because model outputs are
non-deterministic, some observed differences may not represent stable
behavior.

\item \textbf{Single-author evaluation.}
The ground truth, requirement decomposition review, and conflict labels were
produced by the same author. Without an independent reviewer, interpretation
and evaluation bias cannot be ruled out.

\item \textbf{Metric coverage.}
AlignScore is sensitive to wording and requirement decomposition. The
evaluation also does not measure requirements that are missing from the
candidate master prompt.

\end{itemize}

\section{Future Work}
\label{sec:discussion-future}

\begin{itemize}[leftmargin=*]

\item Repeat the evaluation with more runs and more realistic project
histories.

\item Add an evaluation of missing requirements and use more independent
reviewers for ground truth and manual labels.

\item Test the handoff method with more conflicts, compatible additions, and
clear revision prompts.

\item Evaluate whether the governance rules are more effective when the master
prompt is updated after each prompt or when the complete history is
consolidated only once at the end of the session.


\end{itemize}
